%!TEX TS-program = xelatex
\documentclass[11pt]{article}
\usepackage[margin=1in]{geometry}
\usepackage{fontspec}

\newfontfamily\kurdishfont[Script=Arabic]{Amiri}
% For code listings: FreeMono keeps brackets clear and supports Kurdish text.
\newfontfamily\codemonofont{FreeMono}
\newcommand{\codeleftparen}{{\codemonofont(}}
\newcommand{\coderightparen}{{\codemonofont)}}
\newcommand{\codeleftbracket}{{\codemonofont[}}
\newcommand{\coderightbracket}{{\codemonofont]}}
\newcommand{\codeleftbrace}{{\codemonofont\char`\{}}
\newcommand{\coderightbrace}{{\codemonofont\char`\}}}

\usepackage{hyperref}
\usepackage{booktabs}

% Core packages
\usepackage{amsmath,amssymb}
\usepackage{tikz-cd}
\usepackage{multicol}
\usepackage{xcolor}
\usepackage{listings}

% Code listings
\lstdefinestyle{pythoncode}{
  language=Python,
  basicstyle=\footnotesize\codemonofont,
  keywordstyle=\color{blue!70!black},
  commentstyle=\color{green!40!black},
  stringstyle=\color{orange!50!black},
  showstringspaces=false,
  breaklines=true,
  breakatwhitespace=false,
  columns=fullflexible,
  keepspaces=true,
  tabsize=4,
  frame=single,
  numbers=left,
  numberstyle=\scriptsize\codemonofont\color{gray},
  stepnumber=1,
  numbersep=10pt,
  extendedchars=false,
  literate=
    {(}{{\codeleftparen}}1
    {)}{{\coderightparen}}1
    {[}{{\codeleftbracket}}1
    {]}{{\coderightbracket}}1
    {\{}{{\codeleftbrace}}1
    {\}}{{\coderightbrace}}1
}

% Paragraphs
\setlength{\parindent}{0pt}
\setlength{\parskip}{1\baselineskip}

\title{Dataset Processing and Maintenance Toolkit}
\author{Data Engineering Team}
\date{\today}

\begin{document}
\maketitle

\section*{Overview}

This document contains a comprehensive toolkit of Python scripts utilized for cleaning, merging, and maintaining the dataset. Below is a summary table containing links to the code for each step in the pipeline, alongside an explanation of its purpose.

\begin{table}[htbp]
\centering
\small
\renewcommand{\arraystretch}{1.3}
\begin{tabular}{@{} l l p{9cm} @{}}
\toprule
\textbf{Script Name} & \textbf{Section Link} & \textbf{Purpose / Why We Use It} \\
\midrule

\multicolumn{3}{@{}l}{\textbf{Data Processing \& Merging}} \\
\midrule
Dataset Cleaner & \hyperref[sec:cleaner]{Section \ref*{sec:cleaner}} & Removes duplicate entries, strips unnecessary whitespace, and ensures structural integrity of the final JSON data. \\
JSON Collector & \hyperref[sec:collector]{Section \ref*{sec:collector}} & Aggregates scattered finalized JSON files from multiple nested directories into a single consolidated directory, handling collisions automatically. \\
Number Converter & \hyperref[sec:converter]{Section \ref*{sec:converter}} & Standardizes the dataset by transforming Western Arabic numerals (0-9) into Eastern Arabic/Kurdish numerals (0-9) for regional consistency. \\
Dataset Merger & \hyperref[sec:merger]{Section \ref*{sec:merger}} & Combines two major dataset segments into a single file while strictly removing overlaps and missing or malformed records. \\
Batch JSON Merger & \hyperref[sec:batchmerger]{Section \ref*{sec:batchmerger}} & Processes a directory of hundreds of JSON files and consolidates them into one large, unified dataset file while removing irrelevant metadata. \\
\addlinespace

\multicolumn{3}{@{}l}{\textbf{Translation Pipeline}} \\
\midrule
Category Translator & \hyperref[sec:translator]{Section \ref*{sec:translator}} & Leverages external APIs and local rules to dynamically translate categories to standardized English. \\
Advanced Offline Translator & \hyperref[sec:advancedtranslator]{Section \ref*{sec:advancedtranslator}} & Performs advanced word-level decomposition, rule-based extraction, and parallel translation for complex terms. \\
Fallback Translator & \hyperref[sec:fallbacktranslator]{Section \ref*{sec:fallbacktranslator}} & Provides a secondary translation layer targeting only the most resistant non-English categories. \\
\addlinespace

\multicolumn{3}{@{}l}{\textbf{Category Management \& Analysis}} \\
\midrule
Category Sanitizer & \hyperref[sec:sanitizer]{Section \ref*{sec:sanitizer}} & Removes non-ASCII characters, trailing symbols, and normalizes typographic apostrophes. \\
Category Cleaner & \hyperref[sec:catcleaner]{Section \ref*{sec:catcleaner}} & Applies a hard-coded mapping dictionary to definitively resolve stubborn non-English categories. \\
Category Analyzer & \hyperref[sec:analyzer]{Section \ref*{sec:analyzer}} & Scans the dataset for language distribution and provides an overview of category complexity. \\
Category Counter & \hyperref[sec:counter]{Section \ref*{sec:counter}} & A lightweight script to count unique categories and identify non-English entries requiring translation. \\
\addlinespace

\multicolumn{3}{@{}l}{\textbf{Inspection Tools}} \\
\midrule
JSON Inspector & \hyperref[sec:inspector]{Section \ref*{sec:inspector}} & Deeply inspects the JSON structure of the dataset to verify schemas and formatting. \\
Remaining Inspector & \hyperref[sec:remaininspector]{Section \ref*{sec:remaininspector}} & Focuses specifically on inspecting categories that failed automated translations, returning frequencies. \\
\addlinespace

\multicolumn{3}{@{}l}{\textbf{Artifact Maintenance}} \\
\midrule
Artifact Cleaner & \hyperref[sec:artifactcleaner]{Section \ref*{sec:artifactcleaner}} & Cleans and standardizes generated markdown artifacts to ensure translation consistency. \\
Artifact Updater & \hyperref[sec:updater]{Section \ref*{sec:updater}} & A supplementary tool designed to patch the translation markdown artifacts with specific fallback mappings. \\
\bottomrule
\end{tabular}
\caption{Dataset processing scripts categorized by their functional domain.}
\end{table}

\clearpage

\section{Dataset Cleaner} \label{sec:cleaner}
\textbf{Purpose:} We use this script to ensure that the final dataset contains no blank fields or duplicate question-response pairs. It cleans the raw text by normalizing spaces and stripping whitespace, thus dramatically improving dataset quality for downstream tasks.

\begin{lstlisting}[style=pythoncode, caption={dataset\_cleaner.py}]
import json
import re
import os

input_file = 'data/merged_final.json'
output_file = 'data/merged_final_cleaned.json'

def clean_text(text):
    if not isinstance(text, str):
        return ""
    # Remove leading and trailing whitespace
    text = text.strip()
    # Replace multiple spaces with a single space
    text = re.sub(r'\s+', ' ', text)
    return text

def clean_dataset():
    print(f"Loading dataset from {input_file}...")
    with open(input_file, 'r', encoding='utf-8') as f:
        data = json.load(f)
    
    initial_count = len(data)
    print(f"Initial record count: {initial_count}")
    
    cleaned_data = []
    seen = set()
    
    for item in data:
        # Check if question and response exist
        q = item.get('question', '')
        r = item.get('response', '')
        
        # Clean texts
        cleaned_q = clean_text(q)
        cleaned_r = clean_text(r)
        
        # Skip if either is empty
        if not cleaned_q or not cleaned_r:
            continue
            
        # Deduplication based on question and response
        pair_key = (cleaned_q, cleaned_r)
        if pair_key in seen:
            continue
        seen.add(pair_key)
        
        # Update item with cleaned texts
        item['question'] = cleaned_q
        item['response'] = cleaned_r
        
        cleaned_data.append(item)
        
    final_count = len(cleaned_data)
    print(f"Final record count: {final_count}")
    print(f"Removed records: {initial_count - final_count}")
    
    print(f"Saving cleaned dataset to {output_file}...")
    with open(output_file, 'w', encoding='utf-8') as f:
        json.dump(cleaned_data, f, ensure_ascii=False, indent=2)
    print("Done!")

if __name__ == "__main__":
    clean_dataset()
\end{lstlisting}


\section{JSON Collector} \label{sec:collector}
\textbf{Purpose:} In distributed pipelines, outputs often reside in separate directories. This script is used to safely collect all files bearing the \texttt{\_FINAL.json} suffix into one master folder. It performs content comparisons to ensure no data is lost or redundantly copied during naming collisions.

\begin{lstlisting}[style=pythoncode, caption={json\_collector.py}]
import os
import shutil
import filecmp
import sys

def collect_files(source_dir, target_dir):
    # Ensure target directory exists
    os.makedirs(target_dir, exist_ok=True)
    
    # Resolve absolute paths to compare them properly
    source_dir = os.path.abspath(source_dir)
    target_dir = os.path.abspath(target_dir)
    
    copied_count = 0
    collision_count = 0
    skipped_count = 0
    
    print(f"Scanning '{source_dir}' for JSON files ending with '_FINAL.json'...")
    
    for root, dirs, files in os.walk(source_dir):
        # Exclude the target directory from search to avoid double processing
        if os.path.abspath(root).startswith(target_dir):
            continue
            
        for file in files:
            if file.endswith('_FINAL.json'):
                src_path = os.path.join(root, file)
                dest_path = os.path.join(target_dir, file)
                
                already_copied = False
                is_collision = False
                
                if os.path.exists(dest_path):
                    # Check if the existing file has identical content
                    if filecmp.cmp(src_path, dest_path, shallow=False):
                        already_copied = True
                    else:
                        # Naming collision with different content, check numbered files
                        name, ext = os.path.splitext(file)
                        counter = 1
                        while True:
                            check_path = os.path.join(target_dir, f"{name}_{counter}{ext}")
                            if os.path.exists(check_path):
                                if filecmp.cmp(src_path, check_path, shallow=False):
                                    already_copied = True
                                    break
                            else:
                                dest_path = check_path
                                is_collision = True
                                break
                            counter += 1
                
                if already_copied:
                    skipped_count += 1
                    continue
                
                if is_collision:
                    collision_count += 1
                
                # Copy the file
                shutil.copy2(src_path, dest_path)
                print(f"Copied: {src_path} -> {dest_path}")
                copied_count += 1
                
    print("\nSummary:")
    print(f"Total files copied: {copied_count}")
    print(f"Total duplicate files skipped: {skipped_count}")
    if collision_count > 0:
        print(f"Naming collisions resolved: {collision_count}")

if __name__ == '__main__':
    src = sys.argv[1] if len(sys.argv) > 1 else 'data/raw_backups'
    dst = sys.argv[2] if len(sys.argv) > 2 else 'data/all_final_jsons'
    collect_files(src, dst)
\end{lstlisting}


\section{Number Converter} \label{sec:converter}
\textbf{Purpose:} We use this tool to enforce regional typographical standards. It parses all text in the dataset and converts standard Western numerals into Eastern Arabic (Kurdish) digits, ensuring the resulting dataset follows the native language's formatting rules.

\begin{lstlisting}[style=pythoncode, caption={number\_converter.py}]
import json
import os

def convert_to_ckb_numbers(text):
    if not isinstance(text, str):
        return text
        
    kurdish_digits = {
        '0': '٠',
        '1': '١',
        '2': '٢',
        '3': '٣',
        '4': '٤',
        '5': '٥',
        '6': '٦',
        '7': '٧',
        '8': '٨',
        '9': '٩'
    }
    
    result = ""
    for char in text:
        if char in kurdish_digits:
            result += kurdish_digits[char]
        else:
            result += char
    return result

def main():
    file_path = 'data/merged_final.json'
    temp_file_path = 'data/merged_final_temp.json'
    
    print(f"Loading data from {file_path}...")
    with open(file_path, 'r', encoding='utf-8') as f:
        data = json.load(f)
        
    print("Converting numbers...")
    for item in data:
        if 'question' in item:
            item['question'] = convert_to_ckb_numbers(item['question'])
        if 'response' in item:
            item['response'] = convert_to_ckb_numbers(item['response'])
            
    print(f"Saving data to {file_path}...")
    with open(temp_file_path, 'w', encoding='utf-8') as f:
        json.dump(data, f, ensure_ascii=False, indent=2)
        
    os.replace(temp_file_path, file_path)
    print("Done!")

if __name__ == '__main__':
    main()
\end{lstlisting}


\section{Dataset Merger} \label{sec:merger}
\textbf{Purpose:} Sometimes the dataset is processed in separate large chunks. This script merges two independent dataset blocks, strips trailing whitespaces, discards items missing core keys, deduplicates based on Q\&A text, and reassigns sequential IDs.

\begin{lstlisting}[style=pythoncode, caption={dataset\_merger.py}]
import json
import os

def main():
    p1 = 'data/merged_part1.json'
    p2 = 'data/merged_part2.json'
    out_path = 'data/merged_final.json'

    print("Loading dataset 1...")
    with open(p1, 'r', encoding='utf-8') as f:
        d1 = json.load(f)
    print(f"Loaded {len(d1)} items from {p1}")

    print("Loading dataset 2...")
    with open(p2, 'r', encoding='utf-8') as f:
        d2 = json.load(f)
    print(f"Loaded {len(d2)} items from {p2}")

    merged_items = []
    skipped_missing = 0
    skipped_duplicates = 0
    seen_pairs = set()

    for idx, item in enumerate(d1 + d2):
        # Validate keys
        q = item.get('question')
        r = item.get('response')
        cat = item.get('category')

        if not q or not r or not cat:
            skipped_missing += 1
            continue

        q = q.strip()
        r = r.strip()
        cat = cat.strip()

        if not q or not r or not cat:
            skipped_missing += 1
            continue

        # Deduplicate
        pair = (q, r)
        if pair in seen_pairs:
            skipped_duplicates += 1
            continue

        seen_pairs.add(pair)

        # Create cleaned item with only core keys
        cleaned_item = {
            'id': str(len(merged_items) + 1),  # Sequential 1-based string ID
            'category': cat,
            'question': q,
            'response': r
        }
        merged_items.append(cleaned_item)

    print(f"Skipped {skipped_missing} items due to missing/empty question, response, or category.")
    print(f"Skipped {skipped_duplicates} items as exact duplicates (same question and response).")
    print(f"Total merged items: {len(merged_items)}")

    print(f"Writing merged dataset to {out_path}...")
    with open(out_path, 'w', encoding='utf-8') as f:
        json.dump(merged_items, f, ensure_ascii=False, indent=2)

    print("Successfully finished merging!")

if __name__ == '__main__':
    main()
\end{lstlisting}


\section{Batch JSON Merger} \label{sec:batchmerger}
\textbf{Purpose:} After collecting multiple individual JSON files into one folder (via the JSON Collector), this tool loads all of them in a batch process, drops unnecessary document metadata keys, and consolidates all records into one master array.

\begin{lstlisting}[style=pythoncode, caption={batch\_merger.py}]
import os
import json
import sys

def merge_and_clean(source_dir, output_file):
    merged_data = []
    
    # List all JSON files in the directory
    files = sorted([f for f in os.listdir(source_dir) if f.endswith('.json')])
    total_files = len(files)
    
    print(f"Found {total_files} JSON files in '{source_dir}'. Merging...")
    
    for idx, filename in enumerate(files, 1):
        file_path = os.path.join(source_dir, filename)
        try:
            with open(file_path, 'r', encoding='utf-8') as f:
                data = json.load(f)
                
            # Process content, handling lists and individual dictionaries
            if isinstance(data, list):
                for item in data:
                    if isinstance(item, dict):
                        item.pop('document', None)
                    merged_data.append(item)
            elif isinstance(data, dict):
                data.pop('document', None)
                merged_data.append(data)
            else:
                print(f"Warning: Unexpected data structure in {filename}: {type(data)}")
                
        except Exception as e:
            print(f"Error reading {filename}: {e}")
            
        if idx % 500 == 0 or idx == total_files:
            print(f"Processed {idx}/{total_files} files...")
            
    print(f"Writing merged data to {output_file}...")
    with open(output_file, 'w', encoding='utf-8') as f:
        json.dump(merged_data, f, ensure_ascii=False, indent=2)
        
    print(f"Successfully merged {total_files} files into {output_file}.")
    print(f"Total items in merged file: {len(merged_data)}")

if __name__ == '__main__':
    source_dir = sys.argv[1] if len(sys.argv) > 1 else 'data/all_final_jsons'
    output_file = sys.argv[2] if len(sys.argv) > 2 else 'data/merged_final.json'
    merge_and_clean(source_dir, output_file)
\end{lstlisting}


\section{Category Translator} \label{sec:translator}
\textbf{Purpose:} We use this tool to iteratively translate the non-English categories found in the dataset. It detects the language script, uses batching to avoid API limits, and attempts to resolve categories down to unique English terms.

\begin{lstlisting}[style=pythoncode, caption={translate\_categories.py}]
import json
import re
import time
import os
from deep_translator import GoogleTranslator

# Paths
json_path = "data/merged_final.json"
temp_json_path = "data/merged_final_temp.json"
artifact_dir = "artifacts"
artifact_path = os.path.join(artifact_dir, "category_translations.md")

os.makedirs(artifact_dir, exist_ok=True)

# Regex to detect Arabic script
arabic_script_regex = re.compile(r'[\u0600-\u06FF\u0750-\u077F\u08A0-\u08FF]')

def clean_translation(original, translated):
    translated = translated.strip()
    
    # 1. Deduplicate slash terms
    if '/' in translated:
        parts = [p.strip() for p in translated.split('/')]
        parts = [p for p in parts if p]
        seen = []
        for p in parts:
            if p.lower() not in [s.lower() for s in seen]:
                seen.append(p)
        translated = ' / '.join(seen)
        
    # 2. Deduplicate parentheses
    match = re.match(r'^([^(]+)\s*\(([^)]+)\)$', translated)
    if match:
        p1 = match.group(1).strip()
        p2 = match.group(2).strip()
        if p1.lower() == p2.lower():
            translated = p1
            
    # 3. Clean up trailing/leading symbols
    translated = translated.strip(" /()[]{}")
    return translated

print("Loading dataset...")
with open(json_path, 'r', encoding='utf-8') as f:
    data = json.load(f)

# Count frequencies and separate categories
categories_count = {}
for item in data:
    cat = item.get("category")
    if cat is not None:
        categories_count[cat] = categories_count.get(cat, 0) + 1

target_cats = [cat for cat in categories_count.keys() if arabic_script_regex.search(cat)]
print(f"Total unique categories: {len(categories_count)}")
print(f"Categories to translate: {len(target_cats)}")

translator = GoogleTranslator(source='auto', target='en')
translation_mapping = {}
batch_size = 50
total_batches = (len(target_cats) + batch_size - 1) // batch_size

for b_idx in range(total_batches):
    start = b_idx * batch_size
    end = min(start + batch_size, len(target_cats))
    batch = target_cats[start:end]
    
    text = '\n'.join(batch)
    translated_lines = []
    try:
        translated_text = translator.translate(text)
        translated_lines = [line.strip() for line in translated_text.split('\n')]
    except Exception as e:
        print(f"Batch {b_idx+1} translation error. Falling back to individual translation.")
        
    if len(translated_lines) == len(batch):
        for orig, trans in zip(batch, translated_lines):
            translation_mapping[orig] = clean_translation(orig, trans)
    else:
        for item in batch:
            try:
                trans = translator.translate(item)
                translation_mapping[item] = clean_translation(item, trans)
            except Exception:
                translation_mapping[item] = item
                
    time.sleep(0.5)

# Update categories in data
updated_count = 0
for item in data:
    cat = item.get("category")
    if cat in translation_mapping:
        item["category"] = translation_mapping[cat]
        updated_count += 1

print(f"Updated {updated_count} items in the dataset.")

with open(temp_json_path, 'w', encoding='utf-8') as f:
    json.dump(data, f, ensure_ascii=False, indent=2)

os.replace(temp_json_path, json_path)

with open(artifact_path, 'w', encoding='utf-8') as f:
    f.write("# Category Translation Mapping\n\n")
    f.write(f"Total items in dataset: {len(data)}\\\n")
    f.write("| Original Category | Translated Category | Occurrences |\n")
    f.write("| --- | --- | --- |\n")
    sorted_target = sorted(target_cats, key=lambda x: categories_count[x], reverse=True)
    for orig in sorted_target:
        trans = translation_mapping[orig]
        count = categories_count[orig]
        f.write(f"| {orig.replace('|', '\\|')} | {trans.replace('|', '\\|')} | {count} |\n")

print("Artifact written successfully.")
\end{lstlisting}


\section{Artifact Cleaner} \label{sec:artifactcleaner}
\textbf{Purpose:} We use this tool to refine the category translations generated by the translator script. It corrects known discrepancies, merges duplicates, and enforces formatting on the generated markdown artifact to ensure data scientists have an accurate reference sheet.

\begin{lstlisting}[style=pythoncode, caption={clean\_artifact.py}]
import os
import re

artifact_path = "artifacts/category_translations.md"

category_map = {
    'Politics (Siyasət - Politics': 'Politics',
    'Linguistics (Zimanzanî - Linguistics': 'Linguistics',
    'Health (Têndrostî - Health': 'Health',
    'Linguistics / اللغويات': 'Linguistics',
    'Social (Komelayetî': 'Social',
    'Economy (Abûrî - Economy': 'Economy',
    'History / التاریخ': 'History',
    'Zimanşinasî (Linguistics)': 'Linguistics'
}

with open(artifact_path, 'r', encoding='utf-8') as f:
    lines = f.readlines()

new_lines = []
updated_count = 0

for line in lines:
    if line.startswith("|") and not line.startswith("| ---") and not "Original Category" in line:
        parts = [p.strip() for p in line.split("|")]
        if len(parts) >= 4:
            orig = parts[1]
            trans = parts[2]
            count = parts[3]
            
            orig_unescaped = orig.replace('\\|', '|')
            trans_unescaped = trans.replace('\\|', '|')
            
            matched = False
            for k_key, e_val in category_map.items():
                if (k_key in orig_unescaped) or (k_key in trans_unescaped) or (orig_unescaped in k_key):
                    trans_unescaped = e_val
                    matched = True
                    break
            
            english_pattern = re.compile(r"^[a-zA-Z0-9\s\/\(\)\-\.,&:\?!_']*$")
            if not english_pattern.match(trans_unescaped):
                paren_match = re.search(r'\(([^)]+)\)', orig_unescaped)
                if paren_match:
                    trans_unescaped = paren_match.group(1).strip()
                    matched = True
            
            if matched:
                updated_count += 1
                
            orig_escaped = orig_unescaped.replace('|', '\\|')
            trans_escaped = trans_unescaped.replace('|', '\\|')
            new_lines.append(f"| {orig_escaped} | {trans_escaped} | {count} |\n")
        else:
            new_lines.append(line)
    else:
        new_lines.append(line)

with open(artifact_path, 'w', encoding='utf-8') as f:
    f.writelines(new_lines)

print(f"Artifact table updated successfully. Cleaned {updated_count} rows.")
\end{lstlisting}


\section{Category Sanitizer} \label{sec:sanitizer}
\textbf{Purpose:} This script guarantees standard structural integrity for category values across datasets. It enforces standard spaces, replaces typographic quotes, fixes trailing separators, and systematically prevents format divergence.

\begin{lstlisting}[style=pythoncode, caption={sanitize\_categories.py}]
import json
import os
import re

artifact_path = "artifacts/category_translations.md"
json_paths = [
    "data/merged_final.json",
    "data/merged_final_output.json"
]

manual_fixes = {
    'Geography / Géographie': 'Geography',
    'Pოლიtik and Daily Affairs': 'Politics and Daily Affairs',
    'Religious Practice (ʿIbādah)': 'Religious Practice (Ibadah)'
}

def clean_value(val):
    if not val:
        return ""
    val = re.sub(r'\s+', ' ', val).strip()
    if val in manual_fixes:
        val = manual_fixes[val]
    val = val.replace('’', "'").replace('‘', "'").replace('`', "'").replace('´', "'")
    val = val.replace('“', '"').replace('”', '"')
    val = re.sub(r'\s*/\s*', ' / ', val)
    val = re.sub(r'\s*-\s*', ' - ', val)
    return val

for path in json_paths:
    if os.path.exists(path):
        print(f"Cleaning JSON dataset: {path}")
        with open(path, 'r', encoding='utf-8') as f:
            data = json.load(f)
        
        modified_count = 0
        for item in data:
            orig = item.get("category", "")
            cleaned = clean_value(orig)
            if orig != cleaned:
                item["category"] = cleaned
                modified_count += 1
                
        print(f"Standardized {modified_count} records in {os.path.basename(path)}.")
        
        temp_path = path + ".tmp"
        with open(temp_path, 'w', encoding='utf-8') as f:
            json.dump(data, f, ensure_ascii=False, indent=4)
        os.replace(temp_path, path)

if os.path.exists(artifact_path):
    print(f"Cleaning translation report: {artifact_path}")
    with open(artifact_path, "r", encoding="utf-8") as f:
        lines = f.readlines()
        
    new_lines = []
    for line in lines:
        if line.startswith("|") and not line.startswith("| ---") and not "Original Category" in line:
            parts = [p.strip() for p in line.split("|")]
            if len(parts) >= 4:
                orig = parts[1]
                trans = parts[2]
                count = parts[3]
                cleaned_trans = clean_value(trans)
                new_lines.append(f"| {orig} | {cleaned_trans} | {count} |\n")
            else:
                new_lines.append(line)
        else:
            new_lines.append(line)
            
    with open(artifact_path, "w", encoding="utf-8") as f:
        f.writelines(new_lines)
    print("Report file cleaned successfully.")
\end{lstlisting}


\section{Advanced Offline Translator} \label{sec:advancedtranslator}
\textbf{Purpose:} We use this robust translation mechanism for handling complex compound categories from heavily unstandardized texts. It uses offline decomposition, word dictionaries, and high-concurrency API fallbacks. This ensures minimal data loss and maximal standardization on massive multi-source datasets.

\begin{lstlisting}[style=pythoncode, caption={advanced\_translator.py}]
import os
import json
import re
import time
from deep_translator import GoogleTranslator
from concurrent.futures import ThreadPoolExecutor

artifact_path = "artifacts/category_translations.md"
untranslated_json = "data/merged_final_untranslated.json"
output_json = "data/merged_final.json"
temp_output_json = "data/merged_final_temp.json"

def normalize_string(s):
    if not s:
        return ""
    s = s.replace("\u064a", "\u06cc").replace("\u0649", "\u06cc")
    s = s.replace("\u0643", "\u06a9")
    s = s.replace("\u200c", "").replace("\u200b", "")
    s = re.sub(r"\s+", " ", s).strip()
    return s

def clean_term(t):
    return re.sub(r"\(.*?\)", "", normalize_string(t)).strip()

existing_mappings = {}
if os.path.exists(artifact_path):
    with open(artifact_path, "r", encoding="utf-8") as f:
        for line in f:
            if line.startswith("|") and not line.startswith("| ---") and not "Original Category" in line:
                parts = [p.strip() for p in line.split("|")]
                if len(parts) >= 4:
                    orig = parts[1].replace("\\|", "|")
                    trans = parts[2].replace("\\|", "|")
                    existing_mappings[normalize_string(orig)] = trans

roots = {
    "ئابوری": "Economy", "تەندروستی": "Health", "سیاسەت": "Politics", "مێژوو": "History",
    "کۆمەڵایەتی": "Social", "فەرهەنگ": "Culture", "ئەدەب": "Literature", "زانست": "Science",
    "وەرزش": "Sports", "مافی مرۆڤ": "Human Rights", "زمان": "Linguistics"
}
word_mapping = {normalize_string(k): v for k, v in roots.items()}

def is_pure_english(text):
    return bool(re.match(r"^[a-zA-Z0-9\s\/\(\)\-\.,&:\?!_'\''’]*$", text))

def standardize_english_category(cat):
    if not cat: return ""
    cat = re.sub(r"\s+", " ", cat.strip("`'\"“”‘’"))
    words = cat.split()
    capitalized_words = [w.capitalize() if idx == 0 or w.lower() not in {"and","or","of","in","the"} else w.lower() for idx, w in enumerate(words)]
    return " ".join(capitalized_words).strip(" -/&")

with open(untranslated_json, "r", encoding="utf-8") as f:
    data = json.load(f)

unique_cats_counts = {}
for item in data:
    cat = item.get("category", "")
    if cat:
        unique_cats_counts[cat] = unique_cats_counts.get(cat, 0) + 1

resolved_mappings = {}
unresolved_cats = []

for cat in unique_cats_counts.keys():
    cat_norm = normalize_string(cat)
    if is_pure_english(cat):
        resolved_mappings[cat] = standardize_english_category(cat)
    elif cat_norm in existing_mappings:
        resolved_mappings[cat] = standardize_english_category(existing_mappings[cat_norm])
    else:
        unresolved_cats.append(cat)

translator = GoogleTranslator(source="auto", target="en")

def translate_item(orig):
    src = re.sub(r"\(.*?\)", "", orig).strip("<> -/|")
    if not src: src = orig
    for _ in range(3):
        try:
            trans = translator.translate(src)
            return orig, standardize_english_category(trans)
        except Exception:
            time.sleep(1)
    return orig, standardize_english_category(orig)

if unresolved_cats:
    with ThreadPoolExecutor(max_workers=20) as executor:
        results = list(executor.map(translate_item, unresolved_cats))
    for orig, trans in results:
        resolved_mappings[orig] = trans

for item in data:
    cat = item.get("category", "")
    if cat in resolved_mappings:
        item["category"] = resolved_mappings[cat]

with open(temp_output_json, "w", encoding="utf-8") as f:
    json.dump(data, f, ensure_ascii=False, indent=4)
os.replace(temp_output_json, output_json)
\end{lstlisting}


\section{Category Analyzer} \label{sec:analyzer}
\textbf{Purpose:} Analyzes the proportion of English versus non-English categories across the entire corpus, generating a detailed report to help researchers understand the complexity of the dataset's language distribution.

\begin{lstlisting}[style=pythoncode, caption={analyze\_categories.py}]
import json
import re

file_path = "data/merged_final.json"

# Arabic/Kurdish script character range
arabic_script_regex = re.compile(r'[\u0600-\u06FF\u0750-\u077F\u08A0-\u08FF]')

with open(file_path, 'r', encoding='utf-8') as f:
    data = json.load(f)

# Count and unique categories
categories_count = {}
for item in data:
    cat = item.get("category")
    if cat is not None:
        categories_count[cat] = categories_count.get(cat, 0) + 1

kurdish_categories = []
english_categories = []

for cat, count in categories_count.items():
    if arabic_script_regex.search(cat):
        kurdish_categories.append((cat, count))
    else:
        english_categories.append((cat, count))

print(f"Total items in JSON: {len(data)}")
print(f"Total unique categories: {len(categories_count)}")
print(f"Kurdish/mixed categories count: {len(kurdish_categories)}")
print(f"English-only categories count: {len(english_categories)}")

print("\n--- SAMPLE ENGLISH CATEGORIES ---")
for cat, count in sorted(english_categories, key=lambda x: x[1], reverse=True)[:20]:
    print(f"  {cat}: {count}")

print("\n--- ALL KURDISH/MIXED CATEGORIES ---")
for cat, count in sorted(kurdish_categories, key=lambda x: x[1], reverse=True):
    print(f"  {cat}: {count}")
\end{lstlisting}


\section{Category Cleaner} \label{sec:catcleaner}
\textbf{Purpose:} A targeted cleaning script that uses a strictly defined manual mapping dictionary to resolve remaining non-English categories directly in the JSON file without external API calls, ensuring standard formatting.

\begin{lstlisting}[style=pythoncode, caption={clean\_categories\_final.py}]
import json
import os
import re

json_path = "data/merged_final.json"
temp_json_path = "data/merged_final_temp.json"

# Manual mapping dictionary for all remaining non-English categories
category_map = {
    'دین (Religion': 'Religion',
    'دینی (Religion': 'Religion',
    'تاریخ (History': 'History',
    'دینی (Religious': 'Religious',
    'Politics (Siyasət': 'Politics',
    'شعر (Poetry': 'Poetry',
    'Dîn (Religion)': 'Religion',
    'ڕووداو': 'Event',
    'Politics (Siyasət - Politics': 'Politics',
    'زمان (Linguistics': 'Linguistics',
    'قانون (Law': 'Law',
    'Linguistics (Zimananasî': 'Linguistics',
    'بیۆگرافی': 'Biography',
    'گەردوونناسی': 'Astronomy',
    'Fərhəngnigarî (Lexicography)': 'Lexicography',
    'Linguistics (Zimanzanî - Linguistics': 'Linguistics',
    'Linguistics (Zimanşinasi - Linguistics': 'Linguistics',
    'Dînî (Religion)': 'Religion',
    'Dîrok (History)': 'History',
    'Health (Têndrostî - Health': 'Health',
    'Linguistics (Zamanaşinî - Linguistics': 'Linguistics',
    'فەلسەفە / سووفیزم': 'Philosophy / Sufism',
    'تاریخ و فرهنگ': 'History and Culture',
    'تاریخ (Mێژوو': 'History',
    'ئه\u200cده\u200cب و ڕه\u200cخنه': 'Literature and Criticism',
    'Zimanşînasi (Linguistics)': 'Linguistics',
    'ئەخلاق (Ethics': 'Ethics',
    'سیاسەت (Siyasat': 'Politics',
    'مێژوو (Mêjîstû': 'History',
    'سۆفیگەری': 'Sufism',
    'مێژووی ئێران': 'History of Iran',
    'History (Mêjû - History': 'History',
    'Linguistics / اللغويات': 'Linguistics',
    'Felsefe/Zimanşînasi': 'Philosophy / Linguistics',
    'History (Mêjû) / History': 'History',
    'جناح (Cenah': 'Faction',
    'ئەدەب و شیعری': 'Literature and Poetry',
    'اجتماعی': 'Social',
    'History (Mêjû': 'History',
    'Social (Komelayetî': 'Social',
    'Economy (Abûrî': 'Economy',
    'Economy (Abûrî - Economy': 'Economy',
    'History / التاریخ': 'History',
    'Philosophy (Felsəfə': 'Philosophy',
    'Zimanşînasi': 'Linguistics',
    'Zimanşînasi/Felsefe': 'Linguistics / Philosophy',
    'Zimanşinasî (Linguistics)': 'Linguistics',
    'پزیشکی / Medical': 'Medical',
    'سیاسی / Political': 'Political',
    'یاسا و دادوەری (Law and Justice': 'Law and Justice',
    'وەرزش (Sport': 'Sports',
    'ادبی': 'Literature',
    'مێژووی (History': 'History',
    'حەدیس (Hadith': 'Hadith',
    'زمانی (زمان': 'Linguistics',
    'ئاداب و ئەدەب': 'Literature',
    'کۆمه\u200cڵایهتی': 'Social',
    'Adab (Edəb - Literature': 'Literature',
    'Literature (Edəb': 'Literature',
    'Personal Experience (Azmûni Şexsî - Personal Experience': 'Personal Experience',
    'History (Mēzhū - History': 'History',
    'Economy (Abūrī - Economy': 'Economy',
    'Literature (Edêb - Literature': 'Literature',
    'History (Mêjhu - History': 'History',
    'War and War (Ceng u Şer': 'War and War',
    'International Relations (Peywendiyên Nêwdewletiyan': 'International Relations',
    'History and Society (Mêjû u Komelayetî': 'History and Society',
    'Kultūr (Culture)': 'Culture',
    'فەرھەنگی دەستەواژەکان (Idioms and Expressions': 'Idioms and Expressions',
    'Hunər (Art)': 'Art',
    'Colтура / Culture': 'Culture'
}

print("Loading dataset...")
with open(json_path, 'r', encoding='utf-8') as f:
    data = json.load(f)

updated_count = 0
for item in data:
    cat = item.get("category")
    if cat in category_map:
        item["category"] = category_map[cat]
        updated_count += 1

print(f"Successfully updated {updated_count} records in the dataset.")

english_pattern = re.compile(r"^[a-zA-Z0-9\s\/\(\)\-\.,&:\?!_']*$")
remaining_non_english = []
for item in data:
    cat = item.get("category", "")
    if cat and not english_pattern.match(cat):
        remaining_non_english.append(cat)

print(f"Remaining non-English categories after clean: {len(set(remaining_non_english))}")
if remaining_non_english:
    print("Values:", set(remaining_non_english))

with open(temp_json_path, 'w', encoding='utf-8') as f:
    json.dump(data, f, ensure_ascii=False, indent=2)

os.replace(temp_json_path, json_path)
print("Dataset overwritten and updated successfully.")
\end{lstlisting}


\section{Category Counter} \label{sec:counter}
\textbf{Purpose:} A lightweight script to count unique categories and identify non-English entries requiring translation.

\begin{lstlisting}[style=pythoncode, caption={count\_unique\_categories.py}]
import json
import re

file_path = "data/merged_final.json"
arabic_script_regex = re.compile(r'[\u0600-\u06FF\u0750-\u077F\u08A0-\u08FF]')

with open(file_path, 'r', encoding='utf-8') as f:
    data = json.load(f)

unique_cats = set()
kurdish_cats = set()

for item in data:
    cat = item.get("category")
    if cat is not None:
        unique_cats.add(cat)
        if arabic_script_regex.search(cat):
            kurdish_cats.add(cat)

print(f"Total items: {len(data)}")
print(f"Unique categories: {len(unique_cats)}")
print(f"Kurdish categories to translate: {len(kurdish_cats)}")
\end{lstlisting}


\section{JSON Inspector} \label{sec:inspector}
\textbf{Purpose:} Deeply inspects the JSON structure of the dataset to verify schemas and formatting.

\begin{lstlisting}[style=pythoncode, caption={inspect\_json.py}]
import json

file_path = "data/merged_final.json"

with open(file_path, 'r', encoding='utf-8') as f:
    try:
        data = json.load(f)
        print("Type of data:", type(data))
        if isinstance(data, list):
            print("Length of list:", len(data))
            if len(data) > 0:
                print("First item keys:", data[0].keys())
                print("First item sample:", json.dumps(data[0], indent=2, ensure_ascii=False))
                
                categories = set()
                for item in data:
                    cat = item.get("category")
                    if cat is not None:
                        categories.add(str(cat))
                print("\nUnique categories found:", categories)
        elif isinstance(data, dict):
            print("Keys:", list(data.keys()))
            for k in list(data.keys())[:5]:
                print(f"Key: {k}, type of val: {type(data[k])}")
                if isinstance(data[k], dict):
                    print("Val keys:", data[k].keys())
    except Exception as e:
        print("Error reading json:", e)
\end{lstlisting}


\section{Remaining Inspector} \label{sec:remaininspector}
\textbf{Purpose:} Focuses specifically on inspecting categories that failed automated translations, returning frequencies.

\begin{lstlisting}[style=pythoncode, caption={inspect\_remaining.py}]
import json
import re

json_path = "data/merged_final.json"

with open(json_path, 'r', encoding='utf-8') as f:
    data = json.load(f)

cat_counts = {}
for item in data:
    cat = item.get("category", "")
    if cat:
        cat_counts[cat] = cat_counts.get(cat, 0) + 1

english_pattern = re.compile(r'^[a-zA-Z0-9\s\/\(\)\-\.,&:\?!_]*$')

non_english_cats = {}
for cat, count in cat_counts.items():
    if not english_pattern.match(cat):
        non_english_cats[cat] = count

print(f"Total non-English categories found: {len(non_english_cats)}")
print("\nList of categories and counts:")
for cat, count in sorted(non_english_cats.items(), key=lambda x: x[1], reverse=True):
    print(f"  {cat!r}: {count}")
\end{lstlisting}


\section{Fallback Translator} \label{sec:fallbacktranslator}
\textbf{Purpose:} Provides a secondary translation layer targeting only the most resistant non-English categories.

\begin{lstlisting}[style=pythoncode, caption={translate\_remaining.py}]
import json
import re
import time
from deep_translator import GoogleTranslator

json_path = "data/merged_final.json"

with open(json_path, 'r', encoding='utf-8') as f:
    data = json.load(f)

cat_counts = {}
for item in data:
    cat = item.get("category", "")
    if cat:
        cat_counts[cat] = cat_counts.get(cat, 0) + 1

english_pattern = re.compile(r"^[a-zA-Z0-9\s\/\(\)\-\.,&:\?!_']*$")

non_english_cats = []
for cat in cat_counts.keys():
    if not english_pattern.match(cat):
        non_english_cats.append(cat)

print(f"Found {len(non_english_cats)} categories to translate/clean.")

def clean_translation(original, translated):
    translated = translated.strip()
    if '/' in translated:
        parts = [p.strip() for p in translated.split('/')]
        parts = [p for p in parts if p]
        seen = []
        for p in parts:
            if p.lower() not in [s.lower() for s in seen]:
                seen.append(p)
        translated = ' / '.join(seen)
        
    match = re.match(r'^([^(]+)\s*\(([^)]+)\)$', translated)
    if match:
        p1 = match.group(1).strip()
        p2 = match.group(2).strip()
        if p1.lower() == p2.lower():
            translated = p1
            
    translated = translated.strip(" /()[]{}")
    return translated

translator = GoogleTranslator(source='auto', target='en')
mapping = {}

for cat in non_english_cats:
    try:
        translated = translator.translate(cat)
        cleaned = clean_translation(cat, translated)
        mapping[cat] = cleaned
        print(f"{cat!r} -> {cleaned!r}")
        time.sleep(0.2)
    except Exception as e:
        print(f"Error for {cat!r}: {e}")
        mapping[cat] = cat
\end{lstlisting}


\section{Artifact Updater} \label{sec:updater}
\textbf{Purpose:} A supplementary tool designed to patch the translation markdown artifacts with specific fallback mappings.

\begin{lstlisting}[style=pythoncode, caption={update\_artifact.py}]
import os
import re

artifact_path = "artifacts/category_translations.md"

category_map = {
    'دین (Religion': 'Religion',
    'دینی (Religion': 'Religion',
    'تاریخ (History': 'History',
    'دینی (Religious': 'Religious',
    'Politics (Siyasət': 'Politics',
    'شعر (Poetry': 'Poetry',
    'Dîn (Religion)': 'Religion',
    'ڕووداو': 'Event',
    'Politics (Siyasət - Politics': 'Politics',
    'زمان (Linguistics': 'Linguistics',
    'قانون (Law': 'Law',
    'Linguistics (Zimananasî': 'Linguistics',
    'بیۆگرافی': 'Biography',
    'گەردوونناسی': 'Astronomy',
    'Fərhəngnigarî (Lexicography)': 'Lexicography',
    'Linguistics (Zimanzanî - Linguistics': 'Linguistics',
    'Linguistics (Zimanşinasi - Linguistics': 'Linguistics',
    'Dînî (Religion)': 'Religion',
    'Dîrok (History)': 'History',
    'Health (Têndrostî - Health': 'Health',
    'Linguistics (Zamanaşinî - Linguistics': 'Linguistics',
    'فەلسەفە / سووفیزم': 'Philosophy / Sufism',
    'تاریخ و فرهنگ': 'History and Culture',
    'تاریخ (Mێژوو': 'History',
    'ئه\u200cده\u200cب و ڕه\u200cخنه': 'Literature and Criticism',
    'Zimanşînasi (Linguistics)': 'Linguistics',
    'ئەخلاق (Ethics': 'Ethics',
    'سیاسەت (Siyasat': 'Politics',
    'مێژوو (Mêjîstû': 'History',
    'سۆفیگەری': 'Sufism',
    'مێژووی ئێران': 'History of Iran',
    'History (Mêjû - History': 'History',
    'Linguistics / اللغويات': 'Linguistics',
    'Felsefe/Zimanşînasi': 'Philosophy / Linguistics',
    'History (Mêjû) / History': 'History',
    'جناح (Cenah': 'Faction',
    'ئەدەب و شیعری': 'Literature and Poetry',
    'اجتماعی': 'Social',
    'History (Mêjû': 'History',
    'Social (Komelayetî': 'Social',
    'Economy (Abûrî': 'Economy',
    'Economy (Abûrî - Economy': 'Economy',
    'History / التاریخ': 'History',
    'Philosophy (Felsəfə': 'Philosophy',
    'Zimanşînasi': 'Linguistics',
    'Zimanşînasi/Felsefe': 'Linguistics / Philosophy',
    'Zimanşinasî (Linguistics)': 'Linguistics',
    'پزیشکی / Medical': 'Medical',
    'سیاسی / Political': 'Political',
    'یاسا و دادوەری (Law and Justice': 'Law and Justice',
    'وەرزش (Sport': 'Sports',
    'ادبی': 'Literature',
    'مێژووی (History': 'History',
    'حەدیس (Hadith': 'Hadith',
    'زمانی (زمان': 'Linguistics',
    'ئاداب و ئەدەب': 'Literature',
    'کۆمه\u200cڵایهتی': 'Social',
    'Adab (Edəb - Literature': 'Literature',
    'Literature (Edəb': 'Literature',
    'Personal Experience (Azmûni Şexsî - Personal Experience': 'Personal Experience',
    'History (Mēzhū - History': 'History',
    'Economy (Abūrī - Economy': 'Economy',
    'Literature (Edêb - Literature': 'Literature',
    'History (Mêjhu - History': 'History',
    'War and War (Ceng u Şer': 'War and War',
    'International Relations (Peywendiyên Nêwdewletiyan': 'International Relations',
    'History and Society (Mêjû u Komelayetî': 'History and Society',
    'Kultūr (Culture)': 'Culture',
    'فەرھەنگی دەستەواژەکان (Idioms and Expressions': 'Idioms and Expressions',
    'Hunər (Art)': 'Art',
    'Colтура / Culture': 'Culture'
}

with open(artifact_path, 'r', encoding='utf-8') as f:
    lines = f.readlines()

new_lines = []
for line in lines:
    if line.startswith("|") and not line.startswith("| ---") and not "Original Kurdish Category" in line:
        parts = [p.strip() for p in line.split("|")]
        if len(parts) >= 4:
            orig = parts[1]
            trans = parts[2]
            count = parts[3]
            
            orig_unescaped = orig.replace('\\|', '|')
            if orig_unescaped in category_map:
                trans = category_map[orig_unescaped]
            
            orig_escaped = orig_unescaped.replace('|', '\\|')
            trans_escaped = trans.replace('|', '\\|')
            new_lines.append(f"| {orig_escaped} | {trans_escaped} | {count} |\n")
        else:
            new_lines.append(line)
    else:
        new_lines.append(line)

with open(artifact_path, 'w', encoding='utf-8') as f:
    f.writelines(new_lines)

print("Artifact file updated successfully.")
\end{lstlisting}

\end{document}
